The Method of Fundamental Solutions (MFS) approximates a scattered field by a sum of point sources placed inside the scatterer. It needs no mesh and converges rapidly for smooth bodies, but its accuracy depends on where the sources are placed, and the usual choice of a scaled copy of the surface fails for bodies with edges or tips. We study the placement of the sources for the scattering of electromagnetic waves by perfect conductors in three dimensions, using four shapes: a sphere, a torus, a cube and the NASA almond. The guiding rule is that the sources must enclose the singularities of the analytic continuation of the field. For the sphere and the torus the sources lie on a smaller copy of the surface, and the error reaches $2\times10^{-14}$ and $7\times10^{-12}$; for the torus we identify a lower bound for the radius of the source surface. For a convex conducting polyhedron we prove that the scattered field is exactly the field of electric and magnetic currents on the interior medial surface of the body, which shows where the sources must go and why both electric and magnetic dipoles are needed. On the cube the sources therefore cluster towards the edges, and below their smallest depth they are continued by a singular tail whose exponents follow from the edge and vertex expansions of the field. Where the usual layout stalls at an error of $21\%$, the error then decays root-exponentially in the number of unknowns down to $1.5\times10^{-5}$, and the field agrees with a boundary element solution to $8\times10^{-5}$. For the almond, sources on the mid-plane with a tail at the tip give a monostatic radar cross-section that agrees with a published reference simulation to $0.14$\,dB or better on average at frequencies up to $10.25$\,GHz.
